\documentclass[aspectratio=169,10pt]{beamer}
\usetheme{metropolis}
\usepackage{array,booktabs}
\usepackage[T1]{fontenc}
\newcommand{\nn}{\textsc{NixNan}}
\newcommand{\tx}[1]{\texttt{\small #1}}
\title{Instrumenting tensor-core instructions step by step}
\subtitle{\tx{HMMA.16816.F16}, \tx{HMMA.16816.F32}, \tx{HMMA.16816.F32.BF16}}
\author{Ganesh Gopalakrishnan}
\date{SC26 tutorial --- data from \tx{TOOL\_VERBOSE=1} runs of \tx{rd\_tc}}
\begin{document}
\maketitle

\begin{frame}[fragile]{1. What the GPU runs: three HMMA flavours}
\footnotesize
\tx{rd\_tc} (2D reaction--diffusion, one step = two GEMMs) compiles to these SASS instructions (\tx{cuobjdump -sass}):
\begin{semiverbatim}\scriptsize
fp16acc16  HMMA.16816.F16      R22, R16, R22, R8      \textit{// D, A, B, C}
fp16acc32  HMMA.16816.F32      R16, R8.reuse, R22, R16
bf16acc32  HMMA.16816.F32.BF16 R16, R8.reuse, R22, R16
\end{semiverbatim}
\begin{itemize}
\item Operand order is always \textbf{D, A, B, C}: $D = A\times B + C$ for one 16$\times$8$\times$16 tile per warp.
\item Each kernel holds 16 HMMAs (verified: 16 ``Instrumenting instruction HMMA'' lines per run).
\item The first HMMA of every chain has \tx{RZ} as its C operand: the accumulator starts at zero.
\item \nn{} sees only SASS text and operand structure --- it does not know about \tx{mma\_sync}.
\end{itemize}
\end{frame}

\begin{frame}[fragile]{2. The table that describes an HMMA: \tx{instructions.json}}
\footnotesize
The opcode string (with modifiers) is the key; each operand slot gets a \{type, count\}:
\medskip
\centering
\begin{tabular}{@{}l l l l l@{}}
\toprule
 & D (op 0) & A (op 1) & B (op 2) & C (op 3)\\
\midrule
\tx{HMMA.16816.F16} & f16 $\times$4 & f16 $\times$8 & f16 $\times$4 & f16 $\times$4\\
\tx{HMMA.16816.F32} & f32 $\times$4 & f16 $\times$8 & f16 $\times$4 & f32 $\times$4\\
\tx{HMMA.16816.F32.BF16} & f32 $\times$4 & bf16 $\times$8 & bf16 $\times$4 & f32 $\times$4\\
\bottomrule
\end{tabular}
\medskip

\raggedright
\begin{itemize}
\item \texttt{count} = number of \emph{values}, not registers. f16/bf16 are packed two per 32-bit register.
\item \tx{get\_num\_regs}: f16/bf16 $\to\lceil c/2\rceil$, f32 $\to c$, f64 $\to 2c$.
\item Hence F16: D=2, A=4, B=2, C=2 registers. F32 and F32.BF16: D=4, A=4, B=2, C=4.
\item An opcode with no table entry (e.g.\ \tx{F2FP.PACK}) is skipped: nothing is inserted.
\end{itemize}
\end{frame}

\begin{frame}[fragile]{3. Verbose log: what \nn{} decided (first launch only)}
\scriptsize
\begin{semiverbatim}
Instrumenting instruction HMMA.16816.F16 R22, R16, R22, R8 ;
Instrumenting instruction with 3 registers          \textit{// AFTER call}
Instrumenting operand 0. div0: 0, regs: 2, count: 4
Instrumenting instruction with 11 registers         \textit{// BEFORE call}
Instrumenting operand 1. div0: 0, regs: 4
Instrumenting operand 2. div0: 0, regs: 2
Instrumenting operand 3. div0: 0, regs: 2
\end{semiverbatim}
\vspace{-1mm}
\footnotesize
\begin{tabular}{@{}l c c c@{}}
\toprule
 & AFTER (op 0) & BEFORE (ops 1--3) & log lines identical for\\
\midrule
\tx{HMMA.16816.F16} & 3 & 11 & all 16 HMMAs\\
\tx{HMMA.16816.F32} & 5 & 13 & all 16 HMMAs\\
\tx{HMMA.16816.F32.BF16} & 5 & 13 & all 16 HMMAs\\
\bottomrule
\end{tabular}

\smallskip
The number is the device function's \tx{arg\_count}: $\;$ AFTER $=1+\mathrm{regs}(D)$; BEFORE $=3$ reginfo words $+\ \mathrm{regs}(A,B,C)$.\\
F16: $1+2=3$, $3+4+2+2=11$. $\;$F32: $1+4=5$, $3+4+2+4=13$.
\end{frame}

\begin{frame}[fragile]{4. Compact encoding: one 32-bit \tx{reginfo} per operand}
\footnotesize
Bitfield (real struct from \tx{reginfo.cuh}): count:6, type:4, half\_h0, half\_h1, div0, operand:2, num\_regs:4, f64high, reserved.
Type codes: f16=0, bf16=1, f32=2, f64=3. Words below were printed by a g++ program using that struct, filled in as \nn{} fills it (not dumped from the running tool).
\medskip
\centering\scriptsize
\begin{tabular}{@{}l l l l l@{}}
\toprule
 & D & A & B & C\\
\midrule
\tx{F16} & 0x00010C04 & 0x00022C08 & 0x00014C04 & 0x00016C04\\
\tx{F32} & 0x00020C84 & 0x00022C08 & 0x00014C04 & 0x00026C84\\
\tx{F32.BF16} & 0x00020C84 & 0x00022C48 & 0x00014C44 & 0x00026C84\\
\bottomrule
\end{tabular}
\medskip

\raggedright\footnotesize
Decoding \tx{0x00022C48} (A of F32.BF16): count=8, type=1 (bf16), h0=h1=1, operand=1, num\_regs=4.\\
Why compact: the words travel as \emph{call arguments} on the GPU, one per operand, so the device function can classify with no table lookup.
\end{frame}

\begin{frame}[fragile]{5. Which physical registers are picked up}
\footnotesize
\tx{get\_regs} turns each \texttt{REG} operand into \texttt{num\_regs} consecutive \tx{nvbit\_add\_call\_arg\_reg\_val} calls (R$n$, R$n{+}1$, \dots).
\medskip
\centering
\begin{tabular}{@{}l l l l l l@{}}
\toprule
instruction & D (AFTER) & A & B & C (BEFORE)\\
\midrule
\tx{F16 R22,R16,R22,R8} & R22--R23 & R16--R19 & R22--R23 & R8--R9\\
\tx{F32 R16,R8,R22,R16} & R16--R19 & R8--R11 & R22--R23 & R16--R19\\
\tx{BF16 R16,R8,R22,R16} & R16--R19 & R8--R11 & R22--R23 & R16--R19\\
\bottomrule
\end{tabular}
\medskip

\raggedright
\begin{itemize}
\item \textbf{Why two call sites:} D overlaps B (R22) and, for F32, C (R16). \tx{IPOINT\_BEFORE} reads the inputs before HMMA overwrites them; \tx{IPOINT\_AFTER} reads the result D.
\item First HMMA in each chain has C=\tx{RZ}: \nn{} sends \tx{RZ} (index clamped to 255) in every C slot; it reads as 0 (inferred from the code, not separately tested).
\item \tx{.reuse} on A is an operand-cache hint; it does not change the register read.
\end{itemize}
\end{frame}

\begin{frame}[fragile]{6. On the GPU: \tx{nixnan\_check\_regs} classifies the fragments}
\footnotesize
Every thread of the warp passes its own fragment registers (a tile is spread over 32 lanes).
\begin{enumerate}
\item Walk the args: read a \tx{reginfo}, then \texttt{num\_regs} values with \tx{va\_arg}.
\item Classify per type: f16 $\to$ \tx{half2\_classify} (two halves per register), bf16 $\to$ \tx{bf162\_classify}, f32 $\to$ \tx{float\_classify}. Each yields NaN / Inf / $-$Inf / subnormal bits.
\item OR into \tx{exces[operand]}, then shuffle-OR across the warp (the 4 operand slots are why \tx{OPERANDS=4}).
\item Lowest active lane only: \tx{recorder.record(inst\_id, flags, op)} does an \tx{atomicAdd} on a counter indexed by \tx{id<<7 | op<<5 | flags}.
\item If the old count was 0 (first time this instruction/operand/flag-set is seen), push an \tx{exception\_info} to the channel. Repeats are only counted.
\end{enumerate}
This is why the report says ``detected in operand $k$ of instruction \dots of type f16'': the operand slot and the type from the table travel with the record.
\end{frame}

\begin{frame}[fragile]{7. Result: what the instrumented HMMAs reported}
\scriptsize
\begin{semiverbatim}
[infinity] operand 0  HMMA.16816.F16 R26, R4, R30, R26   f16  \textit{// born in D}
[infinity] operand 3  HMMA.16816.F16 R4, R16.reuse, R8, R4    \textit{// reused as C}
[NaN]      operand 0  HMMA.16816.F16 R22, R16, R22, R8        \textit{// 0 x inf}
\end{semiverbatim}
\footnotesize
\begin{itemize}
\item fp16acc16: the Inf appears in operand 0 (D) of \tx{HMMA.16816.F16} at step 230; NaN follows at step 231 (a single Inf spreads along a whole row and column of $U'$ because $P$ is dense).
\item fp16acc32 / bf16acc32: the Inf is only seen as an \emph{input} (operands 1--3) of the next \tx{HMMA.16816.F32[.BF16]}: it is created by \tx{F2FP.PACK}, which has no table entry.
\item Not visible: tile-internal products. \nn{} inspects the register file around the instruction, not the inside of the tensor core.
\end{itemize}
\end{frame}
\end{document}
